% ======================================================================
% TRB Annual Meeting manuscript draft for PyOpenDRIVE.
% ======================================================================

% Document class selection. Leave \manuscriptclass empty to auto-check for a
% .cls file. Auto mode uses a .cls file matching this manuscript name when
% available, then falls back to trbunofficial.cls.
\newcommand{\manuscriptclass}{}
\newcommand{\selectedmanuscriptclass}{}
\makeatletter
\ifx\manuscriptclass\@empty
  \IfFileExists{\jobname.cls}{%
    \renewcommand{\selectedmanuscriptclass}{\jobname}%
  }{%
    \IfFileExists{trbunofficial.cls}{%
      \renewcommand{\selectedmanuscriptclass}{trbunofficial}%
    }{%
      \PackageError{trb_template}{No document class file found}{%
        Add a .cls file that matches this manuscript file name, add
        trbunofficial.cls, or set \string\manuscriptclass\space to the class
        name explicitly.%
      }%
    }%
  }%
\else
  \renewcommand{\selectedmanuscriptclass}{\manuscriptclass}%
\fi
\makeatother

\documentclass[10pt,numbered]{\selectedmanuscriptclass}

\begin{document}

% ---------- Title Page ----------
\title{OpenDRIVE-Based SUMO-CARLA Co-Simulation for Mobility and Energy Assessment in Urban Transportation Networks}

% Edit the affiliation and address before final submission if a formal
% institutional affiliation should be used instead of independent researcher.
\author*{Xiangyong Luo}{Researcher}{Independent Researcher}{luoxiangyong01@gmail.com}[United States]

\submissiondate{May 31, 2026}
\totalpages{}

\maketitle

% ---------- Structured Abstract ----------
% Required headings, in order: Objectives, Methods, Findings, Novelty,
% Practical Applications. Keep the full structured abstract to 300 words or
% fewer. Do not include figures, tables, equations, or undefined acronyms.
\section*{Abstract}

\begingroup
\setlength{\parindent}{0pt}

\noindent\textbf{Objectives:} Urban traffic operations and digital-twin studies
need consistent road networks across microscopic traffic simulation and
high-fidelity driving simulation. This study develops an OpenDRIVE-centered
workflow using PyOpenDRIVE, Simulation of Urban MObility (SUMO), and the CARLA
driving simulator to support reproducible mobility and vehicle energy
assessment from a shared lane-level network.

\par\numberedblankline
\noindent\textbf{Methods:} PyOpenDRIVE is used to inspect and edit an
OpenDRIVE road network, write a reproducible project directory, and generate
SUMO-CARLA co-simulation inputs from a shared \texttt{.xodr} file. The workflow
creates SUMO network, route, vehicle-type, configuration, and run files; CARLA
scripts for loading the same OpenDRIVE network and collecting vehicle telemetry;
and optional SUMO-CARLA synchronization scripts. The case study uses the
repository's Tempe, Arizona, OpenDRIVE and SUMO files and evaluates a 10-minute
smoke-run window. Outputs are analyzed with streaming XML and
comma-separated-value parsers that write reproducible tables, metadata, and
figures.

\par\numberedblankline
\noindent\textbf{Findings:} The implemented workflow prepared a Tempe
OpenDRIVE network containing 5,885 roads, 755 junctions, 15,496 lanes, 4,149
signals, and 243 controllers. A SUMO-only validation run completed two selected
vehicles and produced reproducible mobility summaries, command logs, and
high-resolution figures. CARLA and co-simulation results were not statistically
tested in this smoke run and are explicitly reported as unavailable when the
external CARLA runtime is not executed.

\par\numberedblankline
\noindent\textbf{Novelty:} The study contributes a reproducible OpenDRIVE-based
workflow that links 3D real-world network editing, lane-level preparation, SUMO
mobility-energy simulation setup, CARLA telemetry setup, and
transportation-ready result exports.

\par\numberedblankline
\noindent\textbf{Practical Applications:} Transportation agencies and researchers
can use the workflow to prepare consistent digital-twin simulation cases,
document simulator inputs, and screen traffic operations, electrification, and
control scenarios with transparent reproducibility records.
\par\endgroup
\newpage

% ---------- Main Text ----------
\section{Introduction}\label{sec:intro}

Urban transportation networks are increasingly evaluated through simulation
before field deployment of signal timing plans, connected-vehicle strategies,
electric-fleet policies, or digital-twin applications. These studies often need
two levels of representation. A network-scale traffic simulator is needed for
vehicle movements, delay, emissions, and energy use over many trips. A
high-fidelity driving simulator is useful for three-dimensional visualization,
vehicle telemetry, and automated-driving experiments. The practical difficulty is
that these tools can require different network formats, different coordinate
assumptions, and different execution scripts. A study can therefore lose
reproducibility when the same road network is edited separately for each
simulator.

OpenDRIVE addresses part of this problem by providing a standardized static road
network description for simulation, including roads, reference lines, lanes,
road markings, objects, and signals \citep{asam2026opendrive}. SUMO provides a
widely used microscopic traffic simulation platform for mobility, emission, and
energy analysis \citep{alvarezlopez2018sumo}. CARLA provides an open urban
driving simulator for high-fidelity autonomous-driving and vehicle telemetry
experiments \citep{dosovitskiy2017carla}. However, an applied transportation
study still needs a repeatable workflow that prepares one OpenDRIVE network,
generates the SUMO and CARLA project files, records the exact commands, and
exports analysis tables that can be reviewed.

This paper presents PyOpenDRIVE as an OpenDRIVE-centered workflow layer for
joint SUMO and CARLA studies. PyOpenDRIVE is a Python package adapted from the
libOpenDRIVE project and extended with OpenDRIVE reading, inspection, editing,
SUMO conversion, browser visualization, and a simulation project builder
\citep{libopendrive2023}. The contribution of this paper is not a new traffic
flow model or a new vehicle energy model. Instead, the contribution is a
reproducible workflow that makes simulator preparation and result processing
auditable for transportation studies.

The study has three objectives. First, it defines an OpenDRIVE-centered project
layout that keeps network files, simulator inputs, commands, outputs, and
analysis artifacts together. Second, it implements SUMO, CARLA, and optional
SUMO-CARLA co-simulation builders that can run in dry-run mode when external
simulators are not installed. Third, it demonstrates the workflow on a Tempe
urban network example and reports the generated artifacts and SUMO-only mobility
outputs.

\section{Background}\label{sec:background}

\subsection{OpenDRIVE for Simulation Network Exchange}

OpenDRIVE stores a road network as an XML-based static description with the
file extensions \texttt{.xodr} and \texttt{.xodrz}. The current ASAM standard
describes its application to road-network description, driving simulation, and
traffic simulation \citep{asam2026opendrive}. The format is useful for this
study because it can represent lane-level network geometry and traffic-control
features in one exchange file. That representation does not remove all
interoperability problems. SUMO and CARLA may still interpret geometry,
junctions, lane links, and signals differently. A reproducible workflow must
therefore preserve the source file, record any prepared copy, and report which
simulator-specific files were generated from it.

\subsection{SUMO and CARLA Roles}

SUMO is used in this workflow as the primary source of traffic mobility,
emission, and electric-vehicle energy outputs. Its network conversion, route
generation, and vehicle-type files make it suitable for scenario studies that
vary demand, fleet mix, speed assumptions, and traffic operations. The workflow
creates SUMO configuration files that can request trip information, emissions,
battery outputs, summary statistics, and floating-car data when those outputs
are needed \citep{eclipse2026sumodocs}.

CARLA is used as the high-fidelity simulation and telemetry layer. The workflow
generates scripts that load the same prepared OpenDRIVE network into CARLA
standalone mode, apply a fixed simulation time step, and collect vehicle
positions, orientations, speeds, accelerations, and jerk values. CARLA is not
treated as the default energy model in this workflow unless a separate validated
energy model is added. The optional co-simulation scripts follow the CARLA
SUMO-synchronization pattern and assume that the CARLA server and helper scripts
are installed outside the Python package \citep{carla2026docs}.

\subsection{Recent OpenDRIVE and SUMO-CARLA Co-Simulation Studies}

Recent literature shows why OpenDRIVE-centered preparation remains important
for transportation simulation. Simulator surveys emphasize CARLA's role as an
open-source, high-fidelity 3D environment and identify co-simulation with
traffic simulators as a key capability for autonomous-driving research
\citep{li2024choose, silva2024realistic}. A recent connected-and-automated
vehicle co-simulation survey similarly frames coupled traffic, vehicle,
perception, and communication models as a response to the limits of
single-platform simulation \citep{ashayer2025cosimulation}. Application studies
have used CARLA-SUMO co-simulation for cooperative platooning and merging
\citep{han2022strategic}, the OpenCDA framework and research ecosystem
\citep{xu2021opencda, xu2023opencdaecosystem}, simulation-to-real CDA integration
\citep{zheng2023opencdaros}, infrastructure-assisted trajectory prediction and
risk assessment \citep{cao2025segment}, roadside LiDAR placement
\citep{chen2024roadside}, camera-based traffic-signal reinforcement learning
\citep{azfar2025traffic}, and performance scaling with mesoscopic traffic
simulation \citep{varga2023optimizing}. Verification work such as TSC2CARLA
also demonstrates growing interest in toolchains that translate abstract
driving scenarios into CARLA executions \citep{borchers2025tsc2carla}. Across
these studies, OpenDRIVE and related high-definition map representations provide
the static road-network foundation, while SUMO and CARLA divide network-scale
mobility from high-fidelity vehicle and sensor simulation. The remaining
practical gap is a reproducible network-preparation layer that can inspect,
edit, export, and document the shared OpenDRIVE source before simulator-specific
files are produced.

\section{Methodology}\label{sec:methodology}

\subsection{3.1 PyOpenDRIVE: The 3D Real-World OpenDRIVE Network Viewer and Editor}

The central method in this study is the development of PyOpenDRIVE as a
3D real-world OpenDRIVE network viewer and editor. Existing online OpenDRIVE
viewers are useful for loading an \texttt{.xodr} file, visualizing the road
model in a browser, and in some cases exporting a static 3D mesh. PyOpenDRIVE
extends that viewer concept into an editable, georeferenced, transportation
network preparation environment. In this sense, PyOpenDRIVE is positioned as a
state-of-the-art OpenDRIVE web viewer/editor for this workflow: it combines
OpenDRIVE parsing, lane-level visualization, real-world basemap alignment,
interactive network editing, and \texttt{.xodr}-based SUMO-CARLA
co-simulation environment generation in one reproducible toolchain.

The viewer/editor is implemented with a local Python HTTP server and a MapLibre
GL front end. The server reads an OpenDRIVE file through PyOpenDRIVE, exports
roads, lane polygons, traffic signals, posts, mast arms, and other supported
objects as GeoJSON, and converts OpenDRIVE local coordinates to longitude and
latitude with \texttt{convertXY2LonLat}. When the analyst saves edits, the
server applies the reverse conversion with \texttt{convertLonLat2XY} and writes
supported geometry and attribute changes back to OpenDRIVE XML. This design
keeps the browser responsive while keeping the OpenDRIVE parsing and XML update
logic in Python, where it can be tested and reused by the simulation builder.

The editor is designed for real-world network inspection rather than a detached
schematic view. When the input network contains real-world longitude and
latitude information, the map can be displayed on road, satellite, or terrain
basemaps; otherwise, it falls back to a grid-mesh view for local-coordinate
networks. Users can click and drag lanes, signal heads, signal posts, mast arms,
and environment objects, inspect selected attributes in the OpenDriveViewer
panel, edit lane identifiers and predecessor/successor relationships, add or
delete lanes and signal-support objects, undo the most recent movement, and save
an \texttt{edited\_network.xodr} file through the \texttt{/api/save-xodr}
endpoint. The save path is intentionally conservative: it updates supported
road, lane, signal, and object attributes while preserving OpenDRIVE content
that is not handled by the editor.

This viewer/editor makes PyOpenDRIVE more than a file converter. Starting from
one source \texttt{.xodr} file, PyOpenDRIVE can build the working directory for
a SUMO-CARLA co-simulation case: a prepared OpenDRIVE network, SUMO network and
route inputs, vehicle-type and simulation-configuration files, CARLA scripts for
OpenDRIVE world loading and telemetry collection, and synchronization scripts
for CARLA's SUMO bridge. The viewer/editor provides the reviewable preparation
layer where analysts can verify lane-level geometry, traffic-control assets,
object placement, and coordinate alignment before those simulator files are
created. For a TRB-style workflow, the same design supports reproducibility
because the source \texttt{.xodr}, edited \texttt{.xodr}, generated simulator
inputs, run commands, and analysis artifacts can be stored in the same project
directory.

\subsection{Workflow Design}

The workflow starts with one OpenDRIVE \texttt{.xodr} file and creates a
project directory with separate folders for the prepared network, SUMO files,
CARLA files, co-simulation scripts, and analysis outputs. This structure turns
the \texttt{.xodr} network into an executable SUMO-CARLA co-simulation
environment when the external SUMO and CARLA dependencies are available, and
into a fully auditable dry-run environment when they are not.
Table~\ref{tab:workflow} summarizes the main components. Each build writes a
\texttt{project\_config.json} file containing the source network path, scenario
configuration, build flags, and dry-run status. This file is the main
reproducibility record for the generated case.

\begin{table}[htbp]
  \caption{Workflow Components and Reproducibility Files}
  \label{tab:workflow}
  \centering
  \begin{tabular}{p{0.23\linewidth}p{0.33\linewidth}p{0.34\linewidth}}
    \hline
    Component & Main function & Reproducibility files \\
    \hline
    OpenDRIVE preparation & Copy source \texttt{.xodr}, create prepared
    network, and summarize roads, lanes, signals, and controllers &
    \texttt{network/input.xodr}, \texttt{network/prepared.xodr},
    \texttt{network/network\_summary.json} \\
    SUMO builder & Create or copy SUMO network, route, vehicle-type,
    configuration, and run files & \texttt{sumo/scenario.sumocfg},
    \texttt{sumo/vtypes.add.xml}, \texttt{sumo/build\_sumo\_plan.json} \\
    CARLA builder & Write scripts to load the prepared OpenDRIVE network and
    collect vehicle telemetry & \texttt{carla/load\_opendrive\_world.py},
    \texttt{carla/collect\_telemetry.py} \\
    Co-simulation builder & Write SUMO-CARLA helper commands without starting
    CARLA automatically & \texttt{cosim/run\_synchronization.sh},
    \texttt{cosim/run\_synchronization.bat} \\
    Analysis & Parse SUMO and optional CARLA outputs, then export tables and
    figures & \texttt{scenario\_summary.csv},
    \texttt{mobility\_energy\_table.csv}, and
    \texttt{analysis/figures/} \\
    \hline
  \end{tabular}
\end{table}

\subsection{Network Preparation}

The network preparation step uses standard-library XML parsing to summarize the
OpenDRIVE file without loading all simulator dependencies. The summary counts
roads, junctions, lane sections, lanes, signals, objects, and controllers. It
also records sample road and junction identifiers and the OpenDRIVE header. This
simple summary is not a full semantic validation of OpenDRIVE, but it gives
reviewers and analysts a transparent check that the expected network was used.

\subsection{SUMO Project Generation}

The SUMO builder supports two input modes. If SUMO files already exist beside
the OpenDRIVE file or are named in the scenario configuration, the builder copies
those files into the project. This is useful when a case study already has a
calibrated SUMO network and route file. If those files are not available, the
builder plans or runs \texttt{netconvert} to create a SUMO network from the
prepared OpenDRIVE file and plans or runs \texttt{randomTrips.py} for test
demand generation. Missing SUMO tools are nonfatal in dry-run builds and fatal
only when the user requests a real simulation run.

The builder writes vehicle types for internal-combustion and electric passenger
vehicles. The electric vehicle type includes SUMO battery-device parameters,
while the internal-combustion type uses a configurable emission class. A
scenario file controls electric-vehicle share, maximum speed, driver sigma,
simulation start and end time, and output-file options. When route records are
generated by the workflow, vehicle types are assigned using a seeded random
number generator so the same scenario can be rebuilt reproducibly.

\subsection{CARLA and Co-Simulation Project Generation}

The CARLA builder writes Python scripts instead of directly controlling the
CARLA server during project generation. The loading script connects to the
configured host and port, reads \texttt{network/prepared.xodr}, calls CARLA's
OpenDRIVE world-generation API, and applies fixed-time-step settings. The
telemetry script records frame number, elapsed time, actor identifiers, pose,
velocity, acceleration, speed, acceleration magnitude, and jerk to a CSV file.
Optional ego-vehicle filters can restrict the telemetry file to selected actors.

The co-simulation builder creates shell and Windows batch scripts for CARLA's
SUMO synchronization helpers. The scripts create CARLA-compatible SUMO vehicle
types, convert the OpenDRIVE network for synchronization, and run the
synchronization process with a configurable time step and traffic-light manager.
The generated scripts require \texttt{CARLA\_HOME}; they intentionally do not
start the CARLA server automatically.

\subsection{Analysis}

The analysis module parses SUMO \texttt{tripinfo.xml} with streaming XML parsing
so that large output files can be processed without loading the entire document
tree. Nested \texttt{emissions} and \texttt{battery} records are flattened when
available. The analysis reports completed vehicles, average travel time, route
length, average speed, waiting time, time loss, throughput, total distance, fuel
energy, electricity use, battery energy consumed, regenerated energy, and energy
intensity. CARLA telemetry summaries are written only when a telemetry CSV is
available; otherwise the analysis metadata records why CARLA results were
skipped.

\section{Case Study}\label{sec:case-study}

\subsection{Network and Scenario}

The demonstration uses the Tempe network bundled under
\texttt{paper/examples/datasets/tempe\_net}. The prepared OpenDRIVE file reports
a rectangular header extent of 9,771.55 m east-west and 15,721.36 m north-south.
The case study is intended as a workflow validation example rather than a
calibrated policy evaluation. Table~\ref{tab:network-summary} gives the
OpenDRIVE network counts written by the workflow.

\begin{table}[htbp]
  \caption{Tempe OpenDRIVE Network Summary}
  \label{tab:network-summary}
  \centering
  \begin{tabular}{lr}
    \hline
    Network element & Count \\
    \hline
    Roads & 5,885 \\
    Junctions & 755 \\
    Lane sections & 5,885 \\
    Lanes & 15,496 \\
    Signals & 4,149 \\
    Controllers & 243 \\
    Objects & 0 \\
    \hline
  \end{tabular}
\end{table}

The scenario configuration files include a baseline case, electric-vehicle fleet
share cases, an eco-speed case, and a signal-optimized placeholder case. The
paper demonstration uses the existing SUMO files associated with the Tempe
network and the SUMO-only smoke-run configuration. That configuration runs a
10-minute window from 27,000 s to 27,600 s and disables high-volume emissions,
battery, and floating-car output files to keep the repository example small.

\subsection{Commands}

The dry-run build command creates SUMO, CARLA, and co-simulation files without
requiring a CARLA server:

\begin{verbatim}
python -m pyopendrive.sim.cli build \
  --xodr paper/examples/datasets/tempe_net/tempe.xodr \
  --out paper/examples/sumo_carla_energy/generated/tempe_baseline \
  --scenario paper/examples/sumo_carla_energy/configs/scenario_baseline.yaml \
  --build-sumo --build-carla --build-cosim --dry-run --overwrite
\end{verbatim}

The SUMO-only command builds the project, executes SUMO when it is available,
and analyzes the outputs:

\begin{verbatim}
python -m pyopendrive.sim.cli sumo-only \
  --xodr paper/examples/datasets/tempe_net/tempe.xodr \
  --out paper/examples/sumo_carla_energy/generated/tempe_sumo_only \
  --scenario paper/examples/sumo_carla_energy/configs/scenario_sumo_from_xodr.yaml \
  --overwrite
\end{verbatim}

\section{Results}\label{sec:results}

\subsection{Generated Files}

The workflow generated the expected project folders and reproducibility files.
For the SUMO-only example, the project includes the copied and prepared
OpenDRIVE network, a network summary JSON file, copied SUMO network and route
inputs, a generated SUMO configuration, vehicle-type definitions, a run command
log, and analysis tables and figures. CARLA and co-simulation scripts are
generated by the dry-run build, but CARLA telemetry is not interpreted unless
the CARLA runtime is actually executed.

\subsection{SUMO-Only Mobility Output}

Table~\ref{tab:sumo-results} reports the smoke-run SUMO mobility summary. The
run completed two selected vehicles in the 10-minute example window. Because
detailed energy, emissions, battery, and floating-car outputs were intentionally
disabled for this small repository run, the energy and emissions columns in the
analysis files should be interpreted as unavailable for this demonstration, not
as zero-energy vehicle operation.

\begin{table}[htbp]
  \caption{SUMO-Only Smoke-Run Mobility Metrics}
  \label{tab:sumo-results}
  \centering
  \begin{tabular}{lr}
    \hline
    Metric & Value \\
    \hline
    Completed vehicles & 2 \\
    Average travel time, s & 9.0 \\
    Average route length, m & 148.2 \\
    Average speed, m/s & 16.47 \\
    Average waiting time, s & 0.0 \\
    Average waiting count & 0.0 \\
    Average time loss, s & 0.0 \\
    Throughput, vehicles/h & 720.0 \\
    Total completed distance, km & 0.296 \\
    Detailed energy and emissions & Not analyzed in smoke run \\
    CARLA telemetry & Not executed in SUMO-only run \\
    \hline
  \end{tabular}
\end{table}

\begin{figure}[htbp]
  \centering
  \includegraphics[width=0.88\linewidth]{../examples/sumo_carla_energy/generated/tempe_sumo_only/analysis/figures/mobility_energy_summary.png}
  \caption{Automatically generated mobility-energy summary figure for the
  SUMO-only Tempe smoke run. Energy values are not interpreted for this run
  because detailed energy outputs were disabled in the scenario configuration.}
  \label{fig:mobility-energy-summary}
\end{figure}

Figure~\ref{fig:mobility-energy-summary} shows the generated figure artifact.
The figure is useful as a reproducibility check: the same analysis command writes
the CSV tables and the plotted summary in a deterministic location under the
project's \texttt{analysis/figures} directory. For a final policy-facing study,
the same plotting path can be used after enabling detailed SUMO energy and
emission outputs and running the scenario set across multiple random seeds.

\subsection{CARLA and Co-Simulation Readiness}

The CARLA and co-simulation parts of the workflow were evaluated through file
generation and dependency checks. The generated CARLA files include a script for
loading the prepared OpenDRIVE network into a CARLA world and a script for
recording vehicle telemetry. The generated co-simulation scripts include
commands for CARLA SUMO vehicle-type creation, CARLA-compatible SUMO network
conversion, and synchronization execution. When the CARLA Python API or
\texttt{CARLA\_HOME} is absent, the command-line interface reports these missing
dependencies with explicit messages and does not claim that CARLA results are
available.

\section{Discussion}\label{sec:discussion}

The demonstration shows that PyOpenDRIVE can serve as a practical workflow layer
between an OpenDRIVE network and simulator-specific project files. This is
useful for transportation studies because many reproducibility failures happen
outside the modeling equations: analysts may lose track of which network was
used, which simulator command generated an output, or which external dependency
was needed. The project layout and JSON command logs address those problems in a
simple way that can be reviewed without specialized software.

The current implementation is intentionally conservative about claims. SUMO is
the default source for mobility, emission, and energy measures because its output
files include established trip, emissions, and battery records. CARLA is used
for three-dimensional simulation and telemetry unless an additional validated
energy model is attached. The analysis also avoids fabricating cross-simulator
consistency metrics. If vehicle matching between SUMO and CARLA is unavailable,
the workflow reports the CARLA fields as unavailable rather than estimating
unmatched speed or position differences.

The main limitation of the demonstration is that the reported Tempe run is a
smoke test, not a calibrated multi-scenario evaluation. It verifies that the
network preparation, SUMO execution, parsing, and paper-output generation work
on a real project directory. It does not support conclusions about the energy
effects of electrification, eco-speed operation, or signal optimization. A final
TRB study using this workflow should run the full scenario set, enable detailed
energy and emissions outputs, calibrate demand and vehicle parameters, and use
multiple random seeds.

Another limitation is simulator interpretation. OpenDRIVE improves network
exchange, but it does not guarantee that SUMO and CARLA will represent all
junctions, signals, lanes, and surface geometry identically. The workflow makes
these translations easier to inspect, but final validation should compare
network topology, route feasibility, travel times, and selected trajectory
statistics after both simulators are run.

\section{Conclusions}\label{sec:conclusions}

This paper presents an OpenDRIVE-based PyOpenDRIVE workflow for preparing SUMO,
CARLA, and optional SUMO-CARLA co-simulation projects. The workflow copies and
summarizes a source OpenDRIVE network, writes reproducible project metadata,
generates simulator-specific files and run scripts, and exports analysis tables
and figures. The Tempe example demonstrates that the workflow can prepare a
network with 5,885 roads and 15,496 lanes, run a SUMO-only smoke test, and write
reviewable mobility outputs and figure artifacts.

The workflow is most useful as a reproducible simulation-preparation and
analysis layer for transportation digital-twin studies. Its current results
support workflow feasibility, not policy conclusions. Future work should add
batch scenario execution, calibrated demand inputs, multi-seed statistical
summaries, CARLA trajectory validation, and direct comparison of SUMO and CARLA
vehicle traces where stable vehicle matching is available.

\section{Acknowledgments}

The author thanks the open-source SUMO, CARLA, ASAM OpenDRIVE, and
libOpenDRIVE communities for the tools and standards that make this workflow
possible.

\AIDisclosure{The author used OpenAI ChatGPT 5.5 to polish the manuscrip text and LaTeX formatting from repository files and official author instructions. The author is responsible for verifying the technical claims, references, and final submission metadata before submission.}

\section{Author Contributions}

The author confirms contribution to the paper as follows: study conception and
design: X.~Luo; software implementation and workflow integration: X.~Luo; data
collection and example execution: X.~Luo; analysis and interpretation of
results: X.~Luo; draft manuscript preparation: X.~Luo. The author reviewed the
results and approved the final version of the manuscript.

\section{Declaration of Conflicting Interests}

The author declared no potential conflicts of interest with respect to the
research, authorship, and/or publication of this article.

\section{Funding}

The author disclosed no financial support for the research, authorship, and/or
publication of this article.

\newpage
\begin{thebibliography}{18}
\providecommand{\natexlab}[1]{#1}

\bibitem[{{\'A}lvarez~L{\'o}pez et~al.(2018){\'A}lvarez~L{\'o}pez, Behrisch,
  Bieker-Walz, Erdmann, Fl{\"o}tter{\"o}d, Hilbrich, L{\"u}cken, Rummel,
  Wagner, and Wie{\ss}ner}]{alvarezlopez2018sumo}
{\'A}lvarez~L{\'o}pez, Pablo, Michael Behrisch, Laura Bieker-Walz, Jakob
  Erdmann, Yun-Pang Fl{\"o}tter{\"o}d, Robert Hilbrich, Leonhard L{\"u}cken,
  Johannes Rummel, Peter Wagner, and Evamarie Wie{\ss}ner. 2018. ``Microscopic
  Traffic Simulation Using {SUMO}.'' In \emph{2018 21st International
  Conference on Intelligent Transportation Systems (ITSC)}, 2575--2582. IEEE.
  \url{https://doi.org/10.1109/ITSC.2018.8569938}.

\bibitem[{{ASAM e.V.}(2026)}]{asam2026opendrive}
{ASAM e.V.} 2026. \emph{{ASAM OpenDRIVE 1.9.0}}. Accessed May 31, 2026.
  \url{https://www.asam.net/standards/detail/opendrive/}.

\bibitem[{Ashayer et~al.(2025)Ashayer, Elhag, Sun, and
  Sartipi}]{ashayer2025cosimulation}
Ashayer, Sima, Firas Elhag, Pengyuan Sun, and Mina Sartipi. 2025.
  ``Co-Simulation in Intelligent Transportation Systems: A Survey Targeting
  Connected and Autonomous Vehicles.'' \emph{Transportation Research Record}
  2679 (10). \url{https://doi.org/10.1177/03611981251342248}.

\bibitem[{Azfar et~al.(2025)Azfar, Huang, Tracy, Misiewicz, Liu, and
  Ke}]{azfar2025traffic}
Azfar, Talha, Kaicong Huang, Andrew Tracy, Sandra Misiewicz, Chenxi Liu, and
  Ruimin Ke. 2025. ``Traffic Co-Simulation Framework Empowered by
  Infrastructure Camera Sensing and Reinforcement Learning.'' \emph{Journal of
  Intelligent Transportation Systems}. Published online September 15, 2025.
  \url{https://doi.org/10.1080/15472450.2025.2559410}.

\bibitem[{Borchers et~al.(2025)Borchers, Koopmann, Westhofen, Becker, Putze,
  Grundt, de~Graaff, Kalwa, and Neurohr}]{borchers2025tsc2carla}
Borchers, Philipp, Tjark Koopmann, Lukas Westhofen, Jan Steffen Becker, Lina
  Putze, Dominik Grundt, Thies de~Graaff, Vincent Kalwa, and Christian Neurohr.
  2025. ``TSC2CARLA: An Abstract Scenario-Based Verification Toolchain for
  Automated Driving Systems.'' \emph{Science of Computer Programming} 242:
  103256. \url{https://doi.org/10.1016/j.scico.2024.103256}.

\bibitem[{Cao et~al.(2025)Cao, Shangguan, Visser, Chen, Chai, and
  Cai}]{cao2025segment}
Cao, Yue, Wei Shangguan, Arnoud Visser, Junjie Chen, Linguo Chai, and Baigen
  Cai. 2025. ``Segment-Based Trajectory Prediction and Risk Assessment for
  RSU-Assisted CAVs at Signalized Intersections.'' \emph{IEEE Transactions on
  Intelligent Vehicles} 10 (1): 427--445.
  \url{https://doi.org/10.1109/TIV.2024.3414198}.

\bibitem[{{CARLA Simulator Project}(2026)}]{carla2026docs}
{CARLA Simulator Project}. 2026. \emph{{CARLA} Documentation: OpenDRIVE
  Standalone Mode and SUMO Co-Simulation}. Accessed May 31, 2026.
  \url{https://carla.readthedocs.io/en/latest/}.

\bibitem[{Chen et~al.(2024)Chen, Zheng, and Tan}]{chen2024roadside}
Chen, Yanzhan, Liang Zheng, and Zhen Tan. 2024. ``Roadside LiDAR Placement for
  Cooperative Traffic Detection by a Novel Chance Constrained Stochastic
  Simulation Optimization Approach.'' \emph{Transportation Research Part C:
  Emerging Technologies} 167: 104838.
  \url{https://doi.org/10.1016/j.trc.2024.104838}.

\bibitem[{Dosovitskiy et~al.(2017)Dosovitskiy, Ros, Codevilla, Lopez, and
  Koltun}]{dosovitskiy2017carla}
Dosovitskiy, Alexey, German Ros, Felipe Codevilla, Antonio Lopez, and Vladlen
  Koltun. 2017. ``{CARLA}: An Open Urban Driving Simulator.'' In
  \emph{Proceedings of the 1st Annual Conference on Robot Learning}, 1--16.
  PMLR. \url{https://proceedings.mlr.press/v78/dosovitskiy17a.html}.

\bibitem[{{Eclipse SUMO Project}(2026)}]{eclipse2026sumodocs}
{Eclipse SUMO Project}. 2026. \emph{{SUMO} User Documentation}. Accessed
  May 31, 2026. \url{https://sumo.dlr.de/docs/}.

\bibitem[{Han et~al.(2022)Han, Xu, Xia, Sathyan, Guo, Bujanovic, Leslie, Goli,
  and Ma}]{han2022strategic}
Han, Xu, Runsheng Xu, Xin Xia, Anoop Sathyan, Yi Guo, Pavle Bujanovic, Ed
  Leslie, Mohammad Goli, and Jiaqi Ma. 2022. ``Strategic and Tactical
  Decision-Making for Cooperative Vehicle Platooning with Organized Behavior on
  Multi-Lane Highways.'' \emph{Transportation Research Part C: Emerging
  Technologies} 145: 103952. \url{https://doi.org/10.1016/j.trc.2022.103952}.

\bibitem[{Li et~al.(2024)Li, Yuan, Zhang, Yan, Shen, Wang, and
  Yang}]{li2024choose}
Li, Yueyuan, Wei Yuan, Songan Zhang, Weihao Yan, Qiyuan Shen, Chunxiang Wang,
  and Ming Yang. 2024. ``Choose Your Simulator Wisely: A Review on Open-Source
  Simulators for Autonomous Driving.'' \emph{IEEE Transactions on Intelligent
  Vehicles} 9 (5): 4861--4876.
  \url{https://doi.org/10.1109/TIV.2024.3374044}.

\bibitem[{{libOpenDRIVE Contributors}(2023)}]{libopendrive2023}
{libOpenDRIVE Contributors}. 2023. \emph{{libOpenDRIVE}: Small, Lightweight
  {C++} Library for Handling {OpenDRIVE} Files}. Accessed May 31, 2026.
  \url{https://github.com/pageldev/libOpenDRIVE}.

\bibitem[{Silva et~al.(2024)Silva, Silva, Botelho, and
  Pendao}]{silva2024realistic}
Silva, Ivo, Helder Silva, Fabricio Botelho, and Cristiano Pendao. 2024.
  ``Realistic 3D Simulators for Automotive: A Review of Main Applications and
  Features.'' \emph{Sensors} 24 (18): 5880.
  \url{https://doi.org/10.3390/s24185880}.

\bibitem[{Varga et~al.(2023)Varga, Doba, and Tettamanti}]{varga2023optimizing}
Varga, Balazs, Daniel Doba, and Tamas Tettamanti. 2023. ``Optimizing Vehicle
  Dynamics Co-Simulation Performance by Introducing Mesoscopic Traffic
  Simulation.'' \emph{Simulation Modelling Practice and Theory} 125: 102739.
  \url{https://doi.org/10.1016/j.simpat.2023.102739}.

\bibitem[{Xu et~al.(2021)Xu, Guo, Han, Xia, Xiang, and Ma}]{xu2021opencda}
Xu, Runsheng, Yi Guo, Xu Han, Xin Xia, Hao Xiang, and Jiaqi Ma. 2021.
  ``OpenCDA: An Open Cooperative Driving Automation Framework Integrated with
  Co-Simulation.'' In \emph{2021 IEEE International Intelligent Transportation
  Systems Conference (ITSC)}, 1155--1162. IEEE.
  \url{https://doi.org/10.1109/ITSC48978.2021.9564825}.

\bibitem[{Xu et~al.(2023)Xu, Xiang, Han, Xia, Meng, Chen, Correa-Jullian, and
  Ma}]{xu2023opencdaecosystem}
Xu, Runsheng, Hao Xiang, Xu Han, Xin Xia, Zonglin Meng, Chia-Ju Chen, Camila
  Correa-Jullian, and Jiaqi Ma. 2023. ``The OpenCDA Open-Source Ecosystem for
  Cooperative Driving Automation Research.'' \emph{IEEE Transactions on
  Intelligent Vehicles} 8 (4): 2698--2711.
  \url{https://doi.org/10.1109/TIV.2023.3244948}.

\bibitem[{Zheng et~al.(2023)Zheng, Han, Xia, Gao, Xiang, and
  Ma}]{zheng2023opencdaros}
Zheng, Zhaoliang, Xu Han, Xin Xia, Letian Gao, Hao Xiang, and Jiaqi Ma. 2023.
  ``OpenCDA-ROS: Enabling Seamless Integration of Simulation and Real-World
  Cooperative Driving Automation.'' \emph{IEEE Transactions on Intelligent
  Vehicles} 8 (7): 3775--3780.
  \url{https://doi.org/10.1109/TIV.2023.3298528}.

\end{thebibliography}

\end{document}
